% =====================================================================
%  COURS DE MATHÉMATIQUES — FICHE 1
%  Rappel d'outils et concepts de base en algèbre
%  Document généré en LaTeX avec figures TikZ tracées « from scratch ».
%  Compilation : pdflatex (2 passes pour la table des matières).
% =====================================================================
\documentclass[11pt,a4paper]{article}

% ---------- Encodage & langue ----------
\usepackage[utf8]{inputenc}
\usepackage[T1]{fontenc}
\usepackage{lmodern}
\usepackage[french]{babel}

% ---------- Mathématiques ----------
\usepackage{amsmath,amssymb,mathtools}
\usepackage{icomma}              % virgule décimale correcte en mode maths

% ---------- Unités & nombres ----------
\usepackage{siunitx}
\sisetup{output-decimal-marker={,},per-mode=symbol}

% ---------- Couleurs, boîtes, graphismes ----------
\usepackage{xcolor}
\usepackage{colortbl}   % \rowcolor dans les tableaux
\usepackage{tikz}
\usepackage{pgfplots}
\usepackage[most]{tcolorbox}
\usepackage{enumitem}
\usepackage{array}
\usepackage{booktabs}
\usepackage{multicol}
\usepackage{fancyhdr}
\usepackage{textcomp}            % symbole euro \texteuro
\usepackage[hidelinks]{hyperref}

\usetikzlibrary{arrows.meta,positioning,calc,decorations.pathmorphing,
  decorations.markings,angles,quotes,patterns,backgrounds,
  shapes.geometric,shapes.misc,fadings,intersections,fit}
\pgfplotsset{compat=1.18}

% ---------- Géométrie de page ----------
\usepackage[a4paper,margin=2.2cm,headheight=26pt,headsep=10pt]{geometry}

% =====================================================================
%  PALETTE DE COULEURS  (rose framboise = couleur des maths dans le livre)
% =====================================================================
\definecolor{primary}{RGB}{222,49,99}       % rose framboise
\definecolor{primarydark}{RGB}{150,22,55}
\definecolor{defcol}{RGB}{0,114,178}         % bleu  — définitions
\definecolor{thmcol}{RGB}{0,140,120}         % vert-bleu — théorèmes / propriétés
\definecolor{formulacol}{RGB}{198,93,9}      % orange — formules
\definecolor{remarkcol}{RGB}{120,94,200}     % violet — remarques
\definecolor{attncol}{RGB}{200,40,55}        % rouge — pièges
\definecolor{excol}{RGB}{0,135,90}           % vert — exemples
\definecolor{methcol}{RGB}{52,92,148}        % bleu profond — « comment faire »
\definecolor{curveA}{RGB}{0,90,200}          % courbe bleue
\definecolor{curveB}{RGB}{210,30,40}         % courbe rouge
\definecolor{curveC}{RGB}{0,150,90}          % courbe verte

% =====================================================================
%  STYLE COMMUN DES BOÎTES
% =====================================================================
\tcbset{
  boxbase/.style={enhanced,breakable,boxrule=0.9pt,arc=2.5pt,
     left=3mm,right=3mm,top=2.3mm,bottom=2.3mm,
     fonttitle=\bfseries,coltitle=white,
     toptitle=0.6mm,bottomtitle=0.6mm},
}

% ---- Définition (numérotée) ----
\newtcolorbox[auto counter,number within=section]{definition}[1][]{%
  boxbase,colback=defcol!5,colframe=defcol,
  title={\textsc{Définition}~\thetcbcounter\ \ifx\relax#1\relax\relax\else\textmd{--- \itshape #1}\fi}}

% ---- Théorème / Propriété (numéroté) ----
\newtcolorbox[auto counter,number within=section]{theoreme}[1][]{%
  boxbase,colback=thmcol!6,colframe=thmcol,
  title={\textsc{Propriété / Théorème}~\thetcbcounter\ \ifx\relax#1\relax\relax\else\textmd{--- \itshape #1}\fi}}

% ---- Formule (numérotée) ----
\newtcolorbox[auto counter,number within=section]{formule}[1][]{%
  boxbase,colback=formulacol!6,colframe=formulacol,
  title={\textsc{Formule}~\thetcbcounter\ \ifx\relax#1\relax\relax\else\textmd{--- \itshape #1}\fi}}

% ---- Remarque ----
\newtcolorbox{remarque}[1][]{%
  boxbase,colback=remarkcol!6,colframe=remarkcol,
  title={\textsc{Remarque}\ifx\relax#1\relax\relax\else\textmd{ : \itshape #1}\fi}}

% ---- Attention / piège ----
\newtcolorbox{attention}[1][]{%
  boxbase,colback=attncol!6,colframe=attncol,
  title={\textsc{Attention}\ifx\relax#1\relax\relax\else\textmd{ : \itshape #1}\fi}}

% ---- Exemple ----
\newtcolorbox{exemple}[1][]{%
  boxbase,colback=excol!5,colframe=excol,
  title={\textsc{Exemple}\ifx\relax#1\relax\relax\else\textmd{ : \itshape #1}\fi}}

% ---- Comment faire (méthode pas-à-pas) ----
\newtcolorbox{commentfaire}[1][]{%
  boxbase,colback=methcol!7,colframe=methcol,sharp corners,
  title={\textsc{Comment faire}\ifx\relax#1\relax\relax\else\textmd{ --- \itshape #1}\fi}}

% =====================================================================
%  ENVIRONNEMENT « où : »  (légende des grandeurs d'une formule)
% =====================================================================
\newenvironment{ou}{%
  \par\smallskip\noindent\textbf{où :}\nopagebreak
  \begin{itemize}[leftmargin=1.8em,topsep=2pt,itemsep=1.5pt,parsep=0pt]}%
  {\end{itemize}}

% =====================================================================
%  EN-TÊTES ET PIEDS DE PAGE
% =====================================================================
\pagestyle{fancy}
\fancyhf{}
\fancyhead[L]{\small\textcolor{primarydark}{\textbf{Fiche 1} — Outils de base en algèbre}}
\fancyhead[R]{\small\textcolor{primarydark}{\textsc{Mathématiques}}}
\fancyfoot[C]{\small\thepage}
\renewcommand{\headrulewidth}{0.8pt}
\renewcommand{\headrule}{\hbox to\headwidth{\color{primary}\leaders\hrule height \headrulewidth\hfill}}

% Titres de section en couleur
\usepackage{titlesec}
\titleformat{\section}{\Large\bfseries\color{primarydark}}{\thesection}{0.6em}{}
\titleformat{\subsection}{\large\bfseries\color{primary}}{\thesubsection}{0.6em}{}
\titleformat{\subsubsection}{\normalsize\bfseries\color{primary!80!black}}{\thesubsubsection}{0.6em}{}

% Raccourcis maths
\newcommand{\R}{\mathbb{R}}
\newcommand{\N}{\mathbb{N}}
\newcommand{\Z}{\mathbb{Z}}
\newcommand{\Q}{\mathbb{Q}}
\newcommand{\dd}{\mathrm{d}}
\newcommand{\euro}{\texteuro}

% =====================================================================
\begin{document}
\sloppy

% ---------------------------------------------------------------------
%  PAGE DE TITRE
% ---------------------------------------------------------------------
\begingroup
\thispagestyle{empty}
\centering
\vspace*{1.4cm}

\begin{tikzpicture}
\node[rectangle,rounded corners=8pt,fill=primary,text=white,
      minimum width=15.5cm,minimum height=2.7cm,align=center,inner sep=10pt]
  {{\Huge\bfseries Rappel d'outils}\\[4pt]
   {\LARGE\bfseries et concepts de base en algèbre}};
\end{tikzpicture}

\vspace{0.4cm}
{\large\color{primarydark}\textsc{Cours de mathématiques — Fiche n\textsuperscript{o}\,1}}

\vspace{0.8cm}

% --- Illustration de couverture : exponentielle et logarithme, symétriques par rapport à y=x ---
\begin{tikzpicture}[scale=1.0]
  \begin{axis}[width=10.5cm,height=6.4cm,axis lines=middle,
      xmin=-3.2,xmax=3.6,ymin=-3.2,ymax=3.6,
      xtick=\empty,ytick=\empty,clip=true,
      xlabel={$x$},ylabel={$y$},every axis x label/.style={at={(ticklabel cs:1)},anchor=west},
      every axis y label/.style={at={(ticklabel cs:1)},anchor=south}]
    \addplot[curveA,line width=1.3pt,domain=-3:1.25,samples=80]{exp(x)};
    \addplot[curveB,line width=1.3pt,domain=0.06:3.4,samples=80]{ln(x)};
    \addplot[black!45,dashed,line width=0.8pt,domain=-3.1:3.4]{x};
    \node[curveA] at (axis cs:1.35,3.0){$y=e^{x}$};
    \node[curveB] at (axis cs:3.0,1.45){$y=\ln x$};
    \node[black!55,font=\scriptsize] at (axis cs:2.5,3.0){$y=x$};
  \end{axis}
\end{tikzpicture}

\vfill

\begin{tcolorbox}[boxbase,colback=primary!5,colframe=primary,width=15cm,
    title={\textsc{Objectifs pédagogiques de la fiche}}]
À l'issue de ce cours, l'étudiant·e sera capable de :
\begin{itemize}[leftmargin=1.6em,topsep=2pt,itemsep=1pt]
  \item lire et utiliser les \emph{ensembles de nombres} et leurs notations ($\N_0$, $\Z_0$, $\R_0$, $\R_0^+$) ;
  \item manipuler sans erreur les \emph{puissances}, les \emph{racines} et les \emph{exposants} (entiers, négatifs, fractionnaires) ;
  \item définir et appliquer les \emph{logarithmes} (décimal, népérien, base quelconque) et résoudre des équations exponentielles et logarithmiques ;
  \item \emph{factoriser} une expression (mise en évidence, produits remarquables, méthode de Horner) et résoudre équations et inéquations ;
  \item \emph{résoudre des systèmes} à deux et trois inconnues (substitution, combinaison, pivot de Gauss).
\end{itemize}
\end{tcolorbox}

\vspace{0.5cm}
{\small\color{black!60} Document de synthèse rédigé d'après la Fiche 1 du livre — figures originales tracées en TikZ.}
\par
\endgroup

% ---------------------------------------------------------------------
%  TABLE DES MATIÈRES
% ---------------------------------------------------------------------
\newpage
\renewcommand{\contentsname}{\textcolor{primarydark}{Table des matières}}
\tableofcontents
\thispagestyle{fancy}
\newpage

% =====================================================================
\section{Le langage de l'algèbre : ensembles et notations}
% =====================================================================
Avant tout calcul, il faut savoir \emph{dans quel ensemble} vivent les nombres que l'on manipule.
Une grande partie des erreurs en algèbre vient d'un oubli de domaine : on prend la racine d'un
nombre négatif, on divise par zéro, ou on applique un logarithme à un nombre négatif. Cette
première section fixe le vocabulaire utilisé dans toute la fiche.

\subsection{Les grands ensembles de nombres}

\begin{definition}[les ensembles usuels]
On emboîte les ensembles de nombres les uns dans les autres :
\[ \N \subset \Z \subset \Q \subset \R. \]
\begin{itemize}[leftmargin=1.6em,topsep=2pt,itemsep=1.5pt]
  \item $\N=\{0,1,2,3,\dots\}$ : les \textbf{entiers naturels} ;
  \item $\Z=\{\dots,-2,-1,0,1,2,\dots\}$ : les \textbf{entiers relatifs} (avec les négatifs) ;
  \item $\Q=\bigl\{\tfrac{p}{q}\;:\;p\in\Z,\ q\in\Z,\ q\neq 0\bigr\}$ : les \textbf{rationnels} (fractions) ;
  \item $\R$ : les \textbf{réels} (tous les points d'une droite, y compris $\sqrt2$, $\pi$, $e$\dots).
\end{itemize}
\end{definition}

\begin{center}
\begin{tikzpicture}[scale=1]
  \foreach \r/\c/\lab/\px in {3.4/primary/$\R$/3.05, 2.55/formulacol/$\Q$/2.25, 1.75/thmcol/$\Z$/1.5, 1.0/defcol/$\N$/0.0}{}
  \draw[fill=primary!8,draw=primary,line width=1pt] (0,0) ellipse (3.6 and 2.4);
  \draw[fill=formulacol!10,draw=formulacol,line width=1pt] (-0.4,0) ellipse (2.7 and 1.85);
  \draw[fill=thmcol!12,draw=thmcol,line width=1pt] (-0.9,0) ellipse (1.85 and 1.3);
  \draw[fill=defcol!14,draw=defcol,line width=1pt] (-1.35,0) ellipse (1.0 and 0.78);
  \node[defcol,font=\bfseries] at (-1.35,0) {$\N$};
  \node[thmcol,font=\bfseries] at (-0.1,0) {$\Z$};
  \node[formulacol,font=\bfseries] at (1.25,0) {$\Q$};
  \node[primary,font=\bfseries] at (2.85,0) {$\R$};
  \node[font=\scriptsize,black!70] at (0.55,1.55) {$\sqrt2,\ \pi,\ e$ (irrationnels)};
\end{tikzpicture}
\end{center}

\subsection{Comment lire les notations $\N_0$, $\Z_0$, $\R_0$, $\R_0^+$}

Le livre écrit par exemple : « Soit $a,b\in\R_0$, $m\in\Z_0$, $n\in\N_0$ et $c$ un réel strictement
positif. » Ces indices et exposants \emph{restreignent} l'ensemble. Voici la clé de lecture.

\begin{theoreme}[convention de notation des ensembles privés ou orientés]
\begin{itemize}[leftmargin=1.6em,topsep=2pt,itemsep=1.5pt]
  \item un \textbf{indice $0$} signifie « \emph{privé de} $0$ » : \quad
        $\N_0=\N\setminus\{0\}=\{1,2,3,\dots\}$, \quad $\Z_0=\Z\setminus\{0\}$, \quad $\R_0=\R\setminus\{0\}$ ;
  \item un \textbf{exposant $+$} signifie « \emph{positifs} » (au sens large, $\geq 0$) : \quad
        $\R^{+}=[0,+\infty[$ ;
  \item les deux ensemble, $\R_0^{+}$, désigne les \textbf{réels strictement positifs} : \quad
        $\R_0^{+}=\,]0,+\infty[$.
\end{itemize}
\end{theoreme}

\begin{remarque}
Cette notation est typiquement utilisée dans l'enseignement belge. Retenez la logique :
l'\emph{indice} $0$ enlève le zéro, l'\emph{exposant} de signe choisit le côté de la droite réelle.
Ainsi $\R_0^{+}$ se lit « les réels, sans le zéro, du côté positif » $=\,]0,+\infty[$.
\end{remarque}

\begin{attention}[pourquoi ces restrictions sont vitales]
Les restrictions ne sont pas de la décoration : elles garantissent que les opérations ont un sens.
\begin{itemize}[leftmargin=1.6em,topsep=2pt]
  \item On ne divise jamais par $0$ : un dénominateur doit appartenir à $\R_0$.
  \item Un logarithme $\log_a x$ n'existe que pour $x\in\R_0^{+}$ et une base $a\in\R_0^{+}\setminus\{1\}$.
  \item Une racine carrée réelle $\sqrt{x}$ n'existe que pour $x\in\R^{+}$.
\end{itemize}
\end{attention}

\begin{remarque}[la nature « sans unité » des objets algébriques]
En algèbre pure, les lettres ($a$, $x$, $n$\dots) représentent des \emph{nombres} : ce sont des
grandeurs \textbf{sans dimension}, donc \emph{sans unité physique}. Quand l'algèbre sert à modéliser
une situation concrète (un prix, une longueur), l'unité provient du contexte (l'euro \euro{}, le mètre
\si{\metre}), pas de l'algèbre elle-même. Dans chaque encadré « \textbf{où :} » ci-dessous, nous
précisons donc surtout \emph{l'ensemble d'appartenance} de chaque symbole et ses \emph{contraintes}.
\end{remarque}

% =====================================================================
\section{Puissances, racines et exposants}
% =====================================================================
La \textbf{puissance} est l'outil de calcul le plus fondamental : toute la suite (logarithmes,
factorisation, équations) en dépend. Posons d'abord la définition, puis les règles, puis de
nombreux exemples détaillés.

\begin{definition}[puissance à exposant entier]
Pour un réel $a$ et un entier $n\in\N_0$, la \textbf{puissance} $a^{n}$ est le produit de $n$ facteurs
égaux à $a$ :
\[ a^{n}=\underbrace{a\times a\times\cdots\times a}_{n\ \text{facteurs}}. \]
Le nombre $a$ s'appelle la \textbf{base}, et $n$ l'\textbf{exposant}.
\end{definition}

\subsection{Les règles de calcul des exposants}

\begin{formule}[règles fondamentales des puissances]
Pour $a,b\in\R_0$ et $m,n\in\Z$ :
\[ \boxed{\,a^{n}\times a^{m}=a^{n+m}\,}\quad
   \boxed{\,\dfrac{a^{n}}{a^{m}}=a^{n-m}\,}\quad
   \boxed{\,a^{n}\times b^{n}=(a\times b)^{n}\,}\quad
   \boxed{\,\bigl(a^{n}\bigr)^{m}=a^{nm}\,} \]
\begin{ou}
  \item $a,b$ : les bases, nombres réels non nuls ($a,b\in\R_0$) ;
  \item $n,m$ : les exposants, nombres entiers relatifs ($n,m\in\Z$) ;
  \item le symbole « $\times$ » est le produit ; toutes ces quantités sont des nombres sans unité.
\end{ou}
\end{formule}

\noindent\textbf{Lecture des règles.} Trois idées seulement à mémoriser :
\begin{itemize}[leftmargin=1.6em,topsep=2pt]
  \item \emph{même base, on multiplie} $\Rightarrow$ on \textbf{additionne} les exposants ($a^n\times a^m=a^{n+m}$) ;
  \item \emph{même base, on divise} $\Rightarrow$ on \textbf{soustrait} les exposants ($a^n/a^m=a^{n-m}$) ;
  \item \emph{même exposant, bases différentes} $\Rightarrow$ on peut regrouper les bases ($a^n b^n=(ab)^n$) ;
  \item une \emph{puissance de puissance} $\Rightarrow$ on \textbf{multiplie} les exposants ($(a^n)^m=a^{nm}$).
\end{itemize}

\begin{formule}[exposant nul, négatif et inverse]
\[ \boxed{\,a^{0}=1\,}\ (a\neq 0);\qquad
   \boxed{\,a^{-n}=\dfrac{1}{a^{n}}\,};\qquad
   \boxed{\,a^{-m}\times b^{m}=\bigl(b\,a^{-1}\bigr)^{m}=\Bigl(\dfrac{b}{a}\Bigr)^{m}}. \]
En particulier $a^{-3}=\dfrac{1}{a^{3}}$.
\begin{ou}
  \item $a^{0}=1$ : \emph{toute} base non nulle élevée à la puissance $0$ vaut $1$ ;
  \item $a^{-n}$ : un exposant négatif signifie « \emph{passer à l'inverse} » ($a\in\R_0$) ;
  \item $\Bigl(\dfrac{b}{a}\Bigr)^{m}$ : un quotient élevé à une puissance ($a\in\R_0$).
\end{ou}
\end{formule}

\begin{definition}[racine $n$-ième et exposant fractionnaire]
La \textbf{racine $n$-ième} d'un réel $c>0$ se note $\sqrt[n]{c}$ ; c'est l'unique réel positif dont la
puissance $n$ vaut $c$. On la relie aux puissances par un \textbf{exposant fractionnaire} :
\[ \boxed{\,\sqrt[n]{c}=c^{\frac{1}{n}}\,},\qquad\text{et plus généralement}\qquad
   \sqrt[n]{c^{\,p}}=c^{\frac{p}{n}}. \]
\begin{ou}
  \item $c$ : le radicande, réel strictement positif ($c\in\R_0^{+}$) ;
  \item $n$ : l'indice de la racine, entier $\geq 1$ ($n\in\N_0$) ;
  \item $\dfrac{1}{n}$ ou $\dfrac{p}{n}$ : l'exposant fractionnaire équivalent.
\end{ou}
\end{definition}

\begin{theoreme}[racine d'un carré : valeur absolue]
Attention à un piège fréquent : pour \emph{tout} réel $A$,
\[ \boxed{\,\sqrt{A^{2}}=|A|\,}\qquad\text{(et non simplement } A\text{ !).} \]
La racine carrée renvoie toujours un résultat $\geq 0$ : si $A\geq 0$, $\sqrt{A^2}=A$ ; si $A<0$,
$\sqrt{A^2}=-A>0$. On retiendra cette identité pour la section sur les racines emboîtées.
\end{theoreme}

\subsection{Le signe d'une puissance : exposant pair ou impair}

\begin{theoreme}[règle des signes pour une base négative]
Le signe de $(-a)^{n}$ (avec $a>0$) ne dépend que de la \emph{parité} de l'exposant $n$ :
\[ (-a)^{2n}=+a^{2n}\quad(\text{exposant pair}),\qquad
   (-a)^{2n+1}=-a^{2n+1}\quad(\text{exposant impair}). \]
\end{theoreme}

\noindent Il faut donc distinguer avec soin où s'applique le signe « $-$ ». Détaillons sur l'exemple
$2$, c'est-à-dire $a=2$ :

\begin{exemple}[le « $-$ » dans la base ou devant la puissance ?]
\begin{align*}
  (-2)^{2} &= (-2)\times(-2) = +4 &&\text{(le signe est dans la base : il se neutralise, exposant pair)} \\
  -2^{2} &= -(2\times 2) = -4 &&\text{(le signe est \emph{devant} : il ne participe pas à la puissance)} \\
  -(-2)^{2n} &= -\bigl(2^{2n}\bigr) = -2^{2n} &&\text{(exposant pair $\Rightarrow (-2)^{2n}=2^{2n}$, puis le « $-$ » reste)} \\
  -(-2)^{2n+1} &= -\bigl(-2^{2n+1}\bigr) = +2^{2n+1} &&\text{(exposant impair $\Rightarrow (-2)^{2n+1}=-2^{2n+1}$, le « $-$ » l'annule)}
\end{align*}
\textbf{Conclusion :} $(-2)^{2}\neq -2^{2}$. Le parenthésage change tout.
\end{exemple}

\subsection{Comment faire : simplifier une expression à exposants}

\begin{commentfaire}[ramener une expression sous la forme $c^{k}$]
\begin{enumerate}[leftmargin=2em,topsep=2pt,itemsep=2pt]
  \item \textbf{Convertir} toutes les racines en exposants fractionnaires : $\sqrt[n]{\,\cdot\,}=(\cdot)^{1/n}$.
  \item \textbf{Empiler} les exposants de l'intérieur vers l'extérieur avec $(a^{n})^{m}=a^{nm}$.
  \item \textbf{Regrouper} les exposants : multiplier ceux qui sont emboîtés, additionner ceux d'un même produit.
  \item \textbf{Simplifier} la fraction d'exposants obtenue.
\end{enumerate}
\end{commentfaire}

\begin{exemple}[racines emboîtées — calcul détaillé]
Simplifions l'expression $\displaystyle \sqrt[n]{\sqrt[n]{c^{\,2n^{2}}}}$ pour $c>0$ et $n\in\N_0$.

\medskip
\textbf{Étape 1 — racine intérieure en exposant.} La racine intérieure $\sqrt[n]{c^{2n^2}}$ devient
$\bigl(c^{2n^2}\bigr)^{\frac{1}{n}}$.

\textbf{Étape 2 — racine extérieure en exposant.} Toute l'expression s'écrit alors
\[ \sqrt[n]{\sqrt[n]{c^{\,2n^{2}}}}
   =\left(\Bigl(c^{2n^{2}}\Bigr)^{\frac{1}{n}}\right)^{\!\frac{1}{n}}. \]

\textbf{Étape 3 — on multiplie les exposants emboîtés} (règle $(a^p)^q=a^{pq}$) :
\[ =\Bigl(c^{2n^{2}}\Bigr)^{\frac{1}{n^{2}}}
   = c^{\,2n^{2}\times\frac{1}{n^{2}}}. \]

\textbf{Étape 4 — on simplifie l'exposant :} $\;2n^{2}\times\dfrac{1}{n^{2}}=2$. D'où
\[ \boxed{\;\sqrt[n]{\sqrt[n]{c^{\,2n^{2}}}}=c^{2}.\;} \]
Le résultat ne dépend plus de $n$ : tout le « bruit » se simplifie.
\end{exemple}

\begin{exemple}[manipulations courantes, expliquées une à une]
\begin{align*}
  a^{-m}\times b^{m} &= \frac{b^{m}}{a^{m}} = \Bigl(\frac{b}{a}\Bigr)^{m}
     &&\text{exposant négatif} \to \text{inverse, puis même exposant} \to \text{quotient} \\[2pt]
  \frac{a^{n}}{a^{m}} &= a^{n-m}
     &&\text{même base, division} \to \text{soustraction des exposants} \\[2pt]
  \frac{1}{a^{3}} &= a^{-3}
     &&\text{l'inverse d'une puissance} \to \text{exposant opposé} \\[2pt]
  \bigl(2^{x}\bigr)^{3} &= 2^{3x}
     &&\text{puissance de puissance} \to \text{produit des exposants}
\end{align*}
\end{exemple}

\subsection{La fonction exponentielle $x\mapsto a^{x}$ : croissance et décroissance}

Quand on laisse la base $a$ \emph{fixe} et qu'on fait varier l'exposant $x$, on obtient la fonction
\textbf{exponentielle} de base $a$. Son comportement dépend entièrement de la position de $a$ par
rapport à $1$.

\begin{theoreme}[monotonie de l'exponentielle]
Soit $a\in\R_0^{+}\setminus\{1\}$ et la fonction $f:x\mapsto a^{x}$ définie sur $\R$.
\begin{itemize}[leftmargin=1.6em,topsep=2pt,itemsep=1.5pt]
  \item \textbf{Si $a>1$}, $f$ est \emph{strictement croissante} : un plus grand exposant donne un plus grand résultat.
        \[ a^{x_1}>a^{x_2}\ \text{et}\ a>1 \;\Longrightarrow\; x_1>x_2. \]
  \item \textbf{Si $0<a<1$}, $f$ est \emph{strictement décroissante} : un plus grand exposant donne un plus \emph{petit} résultat.
        \[ a^{x_1}>a^{x_2}\ \text{et}\ 0<a<1 \;\Longrightarrow\; x_1<x_2. \]
  \item Dans les deux cas, $f$ est injective : \quad $a^{x_1}=a^{x_2}\;\Longrightarrow\;x_1=x_2$.
\end{itemize}
\end{theoreme}

\begin{center}
\begin{tikzpicture}
  \begin{axis}[width=11cm,height=6.2cm,axis lines=middle,
      xmin=-3,xmax=3,ymin=-0.4,ymax=5.2,
      xlabel={$x$},ylabel={$a^{x}$},
      xtick={-2,-1,1,2},ytick={1,2,3,4},
      every axis x label/.style={at={(ticklabel cs:1)},anchor=west},
      every axis y label/.style={at={(ticklabel cs:1)},anchor=south},
      legend pos=north west,legend cell align=left,legend style={font=\small,draw=black!30}]
    \addplot[curveA,line width=1.3pt,domain=-3:1.7,samples=80]{exp(x*0.6931)};   % 2^x
    \addplot[curveB,line width=1.3pt,domain=-1.7:3,samples=80]{exp(-x*0.6931)};  % (1/2)^x
    \addplot[black!40,dashed,domain=-3:3]{1};
    \node[curveA] at (axis cs:1.95,4.5){$a=2>1$ (croissante)};
    \node[curveB] at (axis cs:-1.55,4.6){$a=\tfrac12<1$ (décroissante)};
    \node[black!55,font=\scriptsize,anchor=south west] at (axis cs:1.4,1){$a^{0}=1$};
  \end{axis}
\end{tikzpicture}
\end{center}

\begin{remarque}
Les deux courbes passent \emph{toutes} par le point $(0,1)$, puisque $a^{0}=1$ quelle que soit la
base. Remarquez aussi qu'elles sont l'image l'une de l'autre par la symétrie $x\mapsto -x$, car
$\bigl(\tfrac12\bigr)^{x}=2^{-x}$.
\end{remarque}

\subsection{Comment faire : résoudre une équation exponentielle}

\begin{commentfaire}[équation du type $a^{u(x)}=a^{v(x)}$]
L'exponentielle étant injective, l'égalité des puissances équivaut à l'égalité des exposants
(à condition d'avoir \emph{la même base} des deux côtés).
\begin{enumerate}[leftmargin=2em,topsep=2pt,itemsep=2pt]
  \item \textbf{Ramener} les deux membres à une même base $a$ (avec $a>0$, $a\neq 1$).
  \item \textbf{Supprimer} les bases : $a^{u}=a^{v}\Longleftrightarrow u=v$.
  \item \textbf{Résoudre} l'équation ordinaire obtenue sur les exposants.
  \item \textbf{Écrire} l'ensemble solution $S$.
\end{enumerate}
\end{commentfaire}

\begin{exemple}[trois équations exponentielles détaillées]
\textbf{(a)} \;$2^{2x+1}=2^{x}$. Même base $2$, on égale les exposants :
\[ 2x+1=x \;\Longleftrightarrow\; x=-1, \qquad S=\{-1\}. \]

\medskip
\textbf{(b)} \;$2^{2x+1}=2^{-x}$. Toujours la base $2$ :
\[ 2x+1=-x \;\Longleftrightarrow\; 3x=-1 \;\Longleftrightarrow\; x=-\tfrac{1}{3}, \qquad S=\Bigl\{-\tfrac{1}{3}\Bigr\}. \]

\medskip
\textbf{(c)} \;$\displaystyle\Bigl(\frac{2^{2x+1}}{2^{x}}\Bigr)^{3}=1$. On simplifie d'abord l'intérieur
($\tfrac{2^{2x+1}}{2^{x}}=2^{2x+1-x}=2^{x+1}$), puis on monte au cube :
\[ \bigl(2^{x+1}\bigr)^{3}=2^{3x+3}=1=2^{0}
   \;\Longleftrightarrow\; 3x+3=0 \;\Longleftrightarrow\; x=-1,\qquad S=\{-1\}. \]
On a utilisé $1=2^{0}$ pour se ramener à une comparaison d'exposants.
\end{exemple}

% =====================================================================
\section{Les logarithmes}
% =====================================================================
Le logarithme est la \textbf{fonction réciproque} de l'exponentielle : là où l'exponentielle
transforme un exposant en nombre, le logarithme retrouve l'exposant à partir du nombre. C'est
l'outil-clé pour résoudre toute équation où l'inconnue est « en exposant ».

\subsection{Définition et lien avec l'exponentielle}

\begin{definition}[logarithme de base $a$]
Pour une base $a\in\R_0^{+}\setminus\{1\}$ et un nombre $x\in\R_0^{+}$, le \textbf{logarithme de base $a$}
de $x$, noté $\log_a x$, est l'\emph{exposant} qu'il faut donner à $a$ pour obtenir $x$ :
\[ \boxed{\,\log_a x = y \;\Longleftrightarrow\; x = a^{y}\,}. \]
\begin{ou}
  \item $a$ : la base, réel strictement positif et différent de $1$ ($a\in\R_0^{+}\setminus\{1\}$) ;
  \item $x$ : l'argument, réel strictement positif ($x\in\R_0^{+}$) — \emph{on ne prend jamais le log d'un nombre négatif ou nul} ;
  \item $y=\log_a x$ : le résultat, un réel quelconque ($y\in\R$).
\end{ou}
\end{definition}

\begin{center}
\begin{tikzpicture}
  \begin{axis}[width=11cm,height=6cm,axis lines=middle,
      xmin=-2.8,xmax=4.2,ymin=-2.8,ymax=4.2,
      xlabel={$x$},ylabel={$y$},xtick={-2,-1,1,2,3},ytick={-2,-1,1,2,3},
      every axis x label/.style={at={(ticklabel cs:1)},anchor=west},
      every axis y label/.style={at={(ticklabel cs:1)},anchor=south}]
    \addplot[curveA,line width=1.3pt,domain=-2.7:1.4,samples=80]{exp(x)};
    \addplot[curveB,line width=1.3pt,domain=0.06:4,samples=90]{ln(x)};
    \addplot[black!40,dashed,domain=-2.7:4]{x};
    \node[curveA] at (axis cs:1.45,3.5){$y=e^{x}$};
    \node[curveB] at (axis cs:3.4,1.6){$y=\ln x$};
    \node[black!55,font=\scriptsize] at (axis cs:3.2,3.6){$y=x$};
    \fill (axis cs:0,1) circle (1.6pt); \fill (axis cs:1,0) circle (1.6pt);
  \end{axis}
\end{tikzpicture}
\end{center}

\begin{remarque}
Graphiquement, $y=\ln x$ est le \emph{symétrique} de $y=e^{x}$ par rapport à la droite $y=x$ :
c'est la signature de deux fonctions réciproques. Le point $(0,1)$ de l'exponentielle correspond au
point $(1,0)$ du logarithme : $\log_a 1=0$ pour toute base.
\end{remarque}

\subsection{Les propriétés des logarithmes}

\begin{formule}[propriétés fondamentales]
Pour $x,y\in\R_0^{+}$, $a\in\R_0^{+}\setminus\{1\}$ et $n\in\Z_0$ :
\begin{align}
  \log_a(x\times y) &= \log_a x + \log_a y && \text{(produit $\to$ somme)} \tag{L1}\\
  \log_a\!\Bigl(\frac{x}{y}\Bigr) &= \log_a x - \log_a y && \text{(quotient $\to$ différence)} \tag{L2}\\
  \log_a\bigl(x^{n}\bigr) &= n\,\log_a x && \text{(puissance $\to$ produit)} \tag{L3}
\end{align}
\begin{ou}
  \item $x,y$ : arguments strictement positifs ($x,y\in\R_0^{+}$) ;
  \item $a$ : base strictement positive et $\neq 1$ ;
  \item $n$ : exposant entier non nul ($n\in\Z_0$).
\end{ou}
\end{formule}

\noindent\textbf{Idée à retenir.} Le logarithme \emph{abaisse} les opérations d'un cran : il transforme
un \emph{produit} en \emph{somme}, un \emph{quotient} en \emph{différence}, une \emph{puissance} en
\emph{produit}. C'est précisément ce qui permet de « descendre » une inconnue qui se trouve en exposant.

\begin{formule}[valeurs remarquables et changement de base]
\[ \boxed{\,\log_a 1 = 0\,}\quad(\text{car }a^{0}=1);\qquad
   \boxed{\,\log_a a = 1\,}\quad(\text{car }a^{1}=a); \]
\[ \boxed{\,\log_a x = \frac{\log_b x}{\log_b a}\,}\quad\text{(changement de base)};\qquad
   \boxed{\,\log_{a^{n}} x = \frac{1}{n}\,\log_a x\,}. \]
\begin{ou}
  \item $\log_a 1=0$ : le log de $1$ est nul dans \emph{toute} base ;
  \item $\log_a a=1$ : le log de la base elle-même vaut $1$ ;
  \item $b$ : une base auxiliaire quelconque ($b\in\R_0^{+}\setminus\{1\}$) servant à changer de base.
\end{ou}
\end{formule}

\begin{definition}[logarithme décimal et logarithme népérien]
Deux bases sont si fréquentes qu'elles ont une notation propre :
\begin{itemize}[leftmargin=1.6em,topsep=2pt,itemsep=1.5pt]
  \item si $a=10$ : le \textbf{logarithme décimal}, noté $\log x = \log_{10} x$ ;
  \item si $a=e\approx 2{,}718$ : le \textbf{logarithme népérien} (naturel), noté $\ln x=\log_{e} x$.
\end{itemize}
La formule de changement de base les relie : $\displaystyle \log_{10} x=\frac{\ln x}{\ln 10}$, avec $\ln 10\approx 2{,}30$.
\end{definition}

\begin{theoreme}[monotonie du logarithme]
Pour une base $a>1$, $\log_a$ est strictement croissant, donc compatible avec les inégalités :
\[ \log_a x \geq \log_a y \;\Longleftrightarrow\; x\geq y. \]
Pour une base $0<a<1$, $\log_a$ est strictement décroissant : l'inégalité \emph{change de sens}.
\end{theoreme}

\begin{attention}[ne pas confondre $\ln^{2}x$, $\ln x^{2}$ et $\ln(\ln x)$]
Ces écritures sont \emph{différentes} :
\[ \ln^{2}x=(\ln x)^{2}\quad(\text{le log, au carré});\qquad
   \ln x^{2}=2\ln x\quad(\text{log d'un carré, propriété L3}). \]
Ainsi $\ln^{2}x \neq \ln x^{2}$ en général. C'est une source d'erreur classique aux concours.
\end{attention}

\subsection{Exemples de calculs logarithmiques détaillés}

\begin{exemple}[valeurs directes]
\begin{align*}
  \log_a\sqrt[n]{x} &= \log_a x^{1/n} = \tfrac{1}{n}\log_a x &&\text{(L3 avec l'exposant $1/n$)} \\
  \log_a\tfrac{1}{x} &= \log_a x^{-1} = -\log_a x &&\text{(L3 avec l'exposant $-1$)} \\
  \log_{10}10 &= 1, \qquad \ln 1 = 0, \qquad \ln e = 1 &&\text{(valeurs de base)} \\
  \ln e^{2} &= 2\ln e = 2 &&\text{(L3 puis $\ln e=1$)} \\
  \log_{10}10^{2a} &= 2a\,\log_{10}10 = 2a &&\text{(L3)} \\
  \ln 0{,}1 &= \ln\tfrac{1}{10} = -\ln 10 \approx -2{,}30 &&\text{(inverse $\to$ opposé)}
\end{align*}
\end{exemple}

\begin{exemple}[une simplification fine avec les exposants des logarithmes]
Montrons que $\displaystyle \frac{\ln^{2}\!\bigl(2\bigr)}{8}=\frac{1}{2}\,\ln^{2}\!\sqrt{2}$ (résultat utile).
Partons de la définition $\ln^2(2)=(\ln 2)^2$ et de $\ln 2 = \ln\bigl(2^{1/2}\bigr)^{2}=2\ln\sqrt2$, soit
$\ln 2 = 2\ln\sqrt 2$. Alors :
\[ \frac{\ln^{2}(2)}{8}=\frac{(\ln 2)^{2}}{8}
   =\frac{\bigl(2\ln\sqrt2\bigr)^{2}}{8}
   =\frac{4\,(\ln\sqrt2)^{2}}{8}
   =\frac{1}{2}\,\ln^{2}\sqrt2. \]
Chaque étape n'utilise que $\ln 2 = 2\ln\sqrt2$ et l'algèbre des carrés.
\end{exemple}

\subsection{Comment faire : résoudre une équation logarithmique ou « en exposant »}

\begin{commentfaire}[passer du log à l'exponentiel, ou inversement]
\begin{enumerate}[leftmargin=2em,topsep=2pt,itemsep=2pt]
  \item \textbf{Domaine d'abord (C.E.).} Tout argument de logarithme doit être $>0$ : poser les
        \emph{conditions d'existence} avant de calculer.
  \item \textbf{Si l'inconnue est en exposant}, appliquer le logarithme (souvent $\ln$) aux deux membres,
        puis « faire descendre » l'exposant avec $\ln(a^{u})=u\ln a$.
  \item \textbf{Si l'inconnue est dans un logarithme}, repasser à l'exponentiel : $\log_a u = v\Leftrightarrow u=a^{v}$.
  \item \textbf{Résoudre} l'équation obtenue, puis \textbf{vérifier} que les solutions respectent la C.E.
\end{enumerate}
\end{commentfaire}

\begin{exemple}[équation où l'inconnue est en exposant]
Résolvons $e^{2x}=10$. On applique $\ln$ aux deux membres :
\[ \ln e^{2x}=\ln 10 \;\Longrightarrow\; 2x\ln e=\ln 10 \;\Longrightarrow\; 2x=\ln 10
   \;\Longrightarrow\; x=\frac{\ln 10}{2}, \qquad S=\Bigl\{\tfrac{\ln 10}{2}\Bigr\}. \]
On a utilisé $\ln e=1$. Aucune C.E. ici, car $e^{2x}>0$ toujours.
\end{exemple}

\begin{exemple}[équation logarithmique avec condition d'existence]
Résolvons $2\times 10^{5x}=5$.

\textbf{Étape 1 — isoler la puissance :} $\;10^{5x}=\dfrac{5}{2}$.

\textbf{Étape 2 — appliquer $\log_{10}$ :} $\;\log_{10}10^{5x}=\log_{10}\dfrac{5}{2}$, soit $5x=\log_{10}\dfrac52$.

\textbf{Étape 3 — isoler $x$ et développer le quotient (L2) :}
\[ x=\frac{1}{5}\log_{10}\frac{5}{2}=\frac{1}{5}\bigl(\log 5-\log 2\bigr),\qquad
   S=\Bigl\{\tfrac{1}{5}(\log 5-\log 2)\Bigr\}. \]
\end{exemple}

\begin{exemple}[inéquation logarithmique : le domaine est décisif]
Résolvons $\ln(x-1)\geq \ln 2$.

\textbf{C.E. :} l'argument doit être positif : $x-1>0$, donc $x>1$, soit $D_f=\,]1,+\infty[$.

\textbf{Résolution :} $\ln$ étant croissant (base $e>1$), on peut « enlever » les $\ln$ sans changer le sens :
\[ \ln(x-1)\geq\ln 2 \;\Longleftrightarrow\; x-1\geq 2 \;\Longleftrightarrow\; x\geq 3. \]
\textbf{Intersection avec la C.E. :} $\;x\geq 3$ est déjà inclus dans $x>1$. Donc $S=[3,+\infty[$.
\end{exemple}

% =====================================================================
\section{Factorisation et résolution d'équations}
% =====================================================================
\textbf{Factoriser}, c'est récrire une somme comme un \emph{produit} de facteurs. C'est le geste central
de l'algèbre : un produit est nul si et seulement si l'un de ses facteurs est nul, ce qui transforme
une équation compliquée en plusieurs équations simples.

\begin{theoreme}[règle du produit nul]
Pour des réels (ou des expressions) $F$ et $G$ :
\[ \boxed{\,F\times G=0 \;\Longleftrightarrow\; F=0\ \text{ou}\ G=0\,}. \]
C'est pourquoi on cherche toujours à \emph{factoriser} avant de résoudre une équation.
\end{theoreme}

\subsection{Méthode 1 — la mise en évidence}

\begin{definition}[mise en évidence d'un facteur commun]
\textbf{Mettre en évidence}, c'est repérer un \emph{facteur commun} à tous les termes et le « sortir »
devant une parenthèse, en application de la distributivité $k(a+b)=ka+kb$ lue à l'envers.
\end{definition}

\begin{commentfaire}[mettre en évidence]
\begin{enumerate}[leftmargin=2em,topsep=2pt,itemsep=2pt]
  \item \textbf{Repérer} le facteur commun à tous les termes (un nombre, une lettre, une puissance, une parenthèse\dots).
  \item \textbf{Choisir} la plus grande puissance commune et le plus grand diviseur commun des coefficients.
  \item \textbf{Diviser} chaque terme par ce facteur et écrire le reste entre parenthèses.
  \item \textbf{Vérifier} en redéveloppant mentalement.
\end{enumerate}
\end{commentfaire}

\begin{exemple}[trois mises en évidence détaillées]
\textbf{(a)} \;$2x^{2}-3x$. Le facteur commun est $x$ :
\[ 2x^{2}-3x = x\,(2x-3). \]

\textbf{(b)} \;$5\times 2^{2x}-3a\times 2^{2x+1}$. Remarquons $2^{2x+1}=2^{2x}\times 2$, le facteur commun est $2^{2x}$ :
\[ 5\times 2^{2x}-3a\times 2^{2x}\times 2 = 2^{2x}\,(5-3a\times 2) = 2^{2x}\,(5-6a). \]

\textbf{(c)} \;$(x-1)(2x+3)-(x-1)^{2}$. Le facteur commun est la parenthèse $(x-1)$ :
\[ (x-1)\bigl[(2x+3)-(x-1)\bigr] = (x-1)(2x+3-x+1) = (x-1)(x+4). \]
On a soigneusement distribué le signe « $-$ » devant $(x-1)$.
\end{exemple}

\subsection{Méthode 2 — les produits remarquables}

\begin{formule}[les identités remarquables]
\[ \boxed{\,A^{2}-B^{2}=(A-B)(A+B)\,}\qquad\text{(différence de deux carrés)} \]
\[ \boxed{\,A^{2}+2AB+B^{2}=(A+B)^{2}\,}\qquad
   \boxed{\,A^{2}-2AB+B^{2}=(A-B)^{2}\,} \]
\begin{ou}
  \item $A,B$ : deux expressions algébriques quelconques (nombres, monômes, fonctions\dots) ;
  \item le terme central $2AB$ est le « \emph{double produit} » : sa présence (et son signe) signale un carré parfait.
\end{ou}
\end{formule}

\begin{formule}[factorisation du trinôme du second degré]
Pour $A\neq 0$, si le discriminant $\Delta=B^{2}-4AC>0$ :
\[ \boxed{\,Ax^{2}+Bx+C=A\,(x-x_1)(x-x_2)\,},\qquad
   x_{1}=\frac{-B+\sqrt{\Delta}}{2A},\quad x_{2}=\frac{-B-\sqrt{\Delta}}{2A}. \]
Si $\Delta=0$, les deux racines sont confondues : $Ax^{2}+Bx+C=A\,(x-x_1)^{2}$ avec $x_1=\dfrac{-B}{2A}$.
\begin{ou}
  \item $A,B,C$ : les coefficients du trinôme, réels avec $A\neq 0$ ;
  \item $\Delta=B^{2}-4AC$ : le \textbf{discriminant} (sans unité) ; son signe décide du nombre de racines ;
  \item $x_1,x_2$ : les racines (les valeurs qui annulent le trinôme).
\end{ou}
\end{formule}

\begin{center}
\begin{tikzpicture}
  \begin{axis}[width=12.5cm,height=5.6cm,axis lines=middle,
      xmin=-3,xmax=3,ymin=-1.6,ymax=4.4,
      xlabel={$x$},ylabel={$y$},xtick={-2,-1,1,2},ytick={1,2,3},
      every axis x label/.style={at={(ticklabel cs:1)},anchor=west},
      every axis y label/.style={at={(ticklabel cs:1)},anchor=south},
      legend pos=north west,legend cell align=left,legend style={font=\scriptsize,draw=black!30}]
    \addplot[curveC,line width=1.2pt,domain=-2:2,samples=60]{x*x-1};   % Delta>0
    \addplot[curveA,line width=1.2pt,domain=-1.9:1.9,samples=60]{x*x};  % Delta=0
    \addplot[curveB,line width=1.2pt,domain=-1.8:1.8,samples=60]{x*x+1};% Delta<0
    \addplot[black!40,dashed,domain=-3:3]{0};
    \node[curveC,font=\scriptsize] at (axis cs:2.0,3.5){$\Delta>0$ : 2 racines};
    \node[curveA,font=\scriptsize] at (axis cs:-1.55,3.4){$\Delta=0$ : 1 racine double};
    \node[curveB,font=\scriptsize] at (axis cs:0,3.9){$\Delta<0$ : aucune racine réelle};
    \fill[curveC] (axis cs:-1,0) circle (1.6pt); \fill[curveC] (axis cs:1,0) circle (1.6pt);
    \fill[curveA] (axis cs:0,0) circle (1.6pt);
  \end{axis}
\end{tikzpicture}
\end{center}

\begin{exemple}[reconnaître et factoriser, cas par cas]
\textbf{(a) Différence de carrés :} \;$5x^{2}-9=(\sqrt5\,x)^{2}-3^{2}=(\sqrt5\,x-3)(\sqrt5\,x+3)$.

\medskip
\textbf{(b) Carré parfait (signe $-$) :}
\[ 3x^{2}-2\sqrt3\,x+1=(\sqrt3\,x)^{2}-2\times\sqrt3\times 1\times x + 1^{2}=(\sqrt3\,x-1)^{2}. \]
On a identifié $A=\sqrt3\,x$, $B=1$, et vérifié le double produit $2AB=2\sqrt3\,x$.

\medskip
\textbf{(c) Carré parfait (signe $+$) :}
\[ x^{2}+6x+9=x^{2}+2\times 1\times 3\times x+3^{2}=(x+3)^{2}. \]

\medskip
\textbf{(d) Avec mise en évidence préalable :}
\[ 2x^{3}+12x^{2}+18x=2x\,(x^{2}+6x+9)=2x\,(x+3)^{2}. \]

\medskip
\textbf{(e) Avec une exponentielle :}
\[ e^{2x}+2e^{x}+1=(e^{x})^{2}+2\times 1\times e^{x}+1^{2}=(e^{x}+1)^{2}. \]

\medskip
\textbf{(f) Double différence de carrés :}
\[ e^{4x}-1=(e^{2x})^{2}-1^{2}=(e^{2x}-1)(e^{2x}+1)=(e^{x}-1)(e^{x}+1)(e^{2x}+1). \]

\medskip
\textbf{(g) Avec des logarithmes ($x>0$) :}
\[ \ln^{2}x^{2}-3=(\ln x^{2})^{2}-(\sqrt3)^{2}=(\ln x^{2}-\sqrt3)(\ln x^{2}+\sqrt3)=(2\ln x-\sqrt3)(2\ln x+\sqrt3). \]
\end{exemple}

\begin{exemple}[un produit remarquable « malin » sans calculatrice]
Calculons $1\,000\,001^{2}-999\,999^{2}$ \emph{sans} calculatrice, grâce à $A^2-B^2=(A-B)(A+B)$ :
\[ 1\,000\,001^{2}-999\,999^{2}=(1\,000\,001-999\,999)(1\,000\,001+999\,999)=2\times 2\,000\,000=\num{4000000}=4\times 10^{6}. \]
La factorisation évite tout calcul de carré géant.
\end{exemple}

\subsection{Méthode 3 — la factorisation d'un polynôme (méthode de Horner)}

\begin{theoreme}[racine et facteur d'un polynôme]
Un polynôme $P(x)$ est divisible par $(x-a)$ \emph{si et seulement si} $a$ est une racine, c'est-à-dire
$P(a)=0$. On écrit alors
\[ P(x)=(x-a)\times Q(x), \]
où $Q(x)$ est le \textbf{quotient}, un polynôme de degré inférieur d'une unité. Si de plus $Q(b)=0$,
on factorise encore : $P(x)=(x-a)(x-b)\,Q'(x)$.
\end{theoreme}

\noindent La \textbf{méthode de Horner} est un tableau très pratique pour calculer ce quotient $Q(x)$
sans poser de division longue. Illustrons-la sur $P(x)=x^{3}-x^{2}+2x-2$, qui admet $a=1$ comme racine
(car $P(1)=1-1+2-2=0$).

\begin{center}
\begin{tikzpicture}[scale=1,every node/.style={font=\small}]
  % coefficients du haut
  \node at (-1.6,1.2) {$P(x)$ :};
  \node[curveB] at (0,1.2) {$1$};
  \node[curveB] at (2,1.2) {$-1$};
  \node[curveB] at (4,1.2) {$2$};
  \node[curveB] at (6,1.2) {$-2$};
  \node[font=\scriptsize,black!60] at (0,1.75) {$x^{3}$};
  \node[font=\scriptsize,black!60] at (2,1.75) {$x^{2}$};
  \node[font=\scriptsize,black!60] at (4,1.75) {$x^{1}$};
  \node[font=\scriptsize,black!60] at (6,1.75) {$x^{0}$};
  % a = 1 à gauche
  \node[primary,font=\bfseries] at (-1.6,-0.2) {$a=1$};
  % ligne des produits (flèches) et résultats
  \node[curveA] at (2,0.3) {$1$};
  \node[curveA] at (4,0.3) {$0$};
  \node[curveA] at (6,0.3) {$2$};
  % barre horizontale
  \draw[line width=0.9pt] (-0.8,-0.15) -- (6.8,-0.15);
  % résultats (somme)
  \node[primarydark,font=\bfseries] at (0,-0.7) {$1$};
  \node[primarydark,font=\bfseries] at (2,-0.7) {$0$};
  \node[primarydark,font=\bfseries] at (4,-0.7) {$2$};
  \node[excol,font=\bfseries] at (6,-0.7) {$0$};
  \node[excol,font=\scriptsize] at (6.9,-0.7) {reste};
  % flèches descente + multiplication
  \draw[->,curveA,thick] (0,-0.95) to[bend right=35] (1.8,0.2);
  \draw[->,curveA,thick] (2,-0.95) to[bend right=35] (3.8,0.2);
  \draw[->,curveA,thick] (4,-0.95) to[bend right=35] (5.8,0.2);
  \node[curveA,font=\scriptsize] at (1,-1.25) {$\times 1$};
  \node[curveA,font=\scriptsize] at (3,-1.25) {$\times 1$};
  \node[curveA,font=\scriptsize] at (5,-1.25) {$\times 1$};
  % quotient
  \node[anchor=west] at (-1.6,-1.9) {$Q(x)=1\,x^{2}+0\,x+2=x^{2}+2$};
  % résultat encadré
  \node[draw=primary,rounded corners,fill=primary!6,anchor=west,inner sep=5pt] at (2.6,-2.5)
       {$P(x)=(x-1)\times(x^{2}+2)$};
\end{tikzpicture}
\end{center}

\begin{commentfaire}[schéma de Horner, pas à pas]
\begin{enumerate}[leftmargin=2em,topsep=2pt,itemsep=2pt]
  \item \textbf{Aligner} les coefficients de $P$ dans l'ordre des puissances \emph{décroissantes}
        (mettre un $0$ pour toute puissance manquante).
  \item \textbf{Trouver} une racine évidente $a$ (souvent $\pm1,\pm2,\dots$) telle que $P(a)=0$.
  \item \textbf{Abaisser} le premier coefficient tel quel.
  \item \textbf{Multiplier} par $a$, écrire le résultat sous le coefficient suivant, puis \textbf{additionner}.
        Répéter jusqu'au bout.
  \item \textbf{Lire} le résultat : les sommes (sauf la dernière) sont les coefficients de $Q(x)$ ;
        la \emph{dernière} somme est le \textbf{reste} (il doit valoir $0$ si $a$ est bien racine).
\end{enumerate}
\end{commentfaire}

\begin{exemple}[factorisation complète de $4x^{3}-8x^{2}+5x-1$]
Soit $P(x)=4x^{3}-8x^{2}+5x-1$.

\textbf{Étape 1 — racine évidente :} $P(1)=4-8+5-1=0$, donc $(x-1)$ divise $P$.

\textbf{Étape 2 — Horner} donne le quotient $Q(x)=4x^{2}-4x+1$, d'où
\[ P(x)=(x-1)\bigl(4x^{2}-4x+1\bigr). \]

\textbf{Étape 3 — reconnaître un carré parfait :} $4x^{2}-4x+1=(2x)^{2}-2\times 2x\times 1+1^{2}=(2x-1)^{2}$.

\textbf{Conclusion :} $\;P(x)=(x-1)(2x-1)^{2}$.

\textbf{Étape 4 — résoudre $P(x)=0$ (règle du produit nul) :}
\[ (x-1)(2x-1)^{2}=0 \;\Longleftrightarrow\; x-1=0\ \text{ou}\ 2x-1=0
   \;\Longleftrightarrow\; x=1\ \text{ou}\ x=\tfrac12,\qquad S=\Bigl\{\tfrac12,\,1\Bigr\}. \]
\end{exemple}

\subsection{Comment faire : simplifier une fraction rationnelle}

\begin{commentfaire}[simplifier $\dfrac{P(x)}{Q(x)}$]
\begin{enumerate}[leftmargin=2em,topsep=2pt,itemsep=2pt]
  \item \textbf{Conditions d'existence (C.E.) :} chercher les valeurs qui annulent le \emph{dénominateur} et les exclure.
  \item \textbf{Factoriser} entièrement le numérateur \emph{et} le dénominateur.
  \item \textbf{Simplifier} les facteurs communs (autorisé seulement parce qu'ils sont non nuls grâce à la C.E.).
  \item \textbf{Conclure} en rappelant la C.E. (la fraction simplifiée n'est égale à l'originale que là où elle existe).
\end{enumerate}
\end{commentfaire}

\begin{exemple}[simplification détaillée]
Simplifions $\displaystyle \frac{4x^{3}-8x^{2}+5x-1}{2x^{2}-3x+1}$.

\textbf{C.E. :} on annule le dénominateur $2x^{2}-3x+1=0$. Discriminant : $\Delta=(-3)^{2}-4\times2\times1=1$, racines
\[ x_1=\frac{3+\sqrt1}{4}=1,\qquad x_2=\frac{3-\sqrt1}{4}=\frac{1}{2}. \]
Donc C.E. : $x\neq 1$ et $x\neq\tfrac12$, et $2x^{2}-3x+1=2\bigl(x-\tfrac12\bigr)(x-1)=(2x-1)(x-1)$.

\textbf{Numérateur} (vu ci-dessus) : $4x^{3}-8x^{2}+5x-1=(x-1)(2x-1)^{2}$.

\textbf{Simplification :}
\[ \frac{4x^{3}-8x^{2}+5x-1}{2x^{2}-3x+1}
   =\frac{(x-1)(2x-1)^{2}}{(2x-1)(x-1)}
   =2x-1,\qquad \text{si } x\neq 1\ \text{et}\ x\neq\tfrac12. \]
\end{exemple}

\subsection{Comment faire : résoudre une inéquation avec un tableau de signes}

\begin{commentfaire}[tableau de signes d'un quotient ou d'un produit]
\begin{enumerate}[leftmargin=2em,topsep=2pt,itemsep=2pt]
  \item \textbf{C.E. :} exclure les valeurs annulant un dénominateur.
  \item \textbf{Factoriser} pour obtenir un produit/quotient de facteurs du \emph{premier degré}.
  \item \textbf{Valeurs critiques :} repérer les zéros de chaque facteur ; les placer dans l'ordre croissant sur la ligne du haut.
  \item \textbf{Signe de chaque facteur} sur chaque intervalle (un facteur $x-r$ est négatif avant $r$, positif après).
  \item \textbf{Multiplier les signes} colonne par colonne (règle des signes), en marquant $0$ aux zéros et « $\nexists$ » aux valeurs interdites.
  \item \textbf{Lire} l'ensemble solution selon le sens de l'inéquation ($>0$, $\geq 0$, $<0$\dots).
\end{enumerate}
\end{commentfaire}

\begin{exemple}[résolution complète d'une inéquation]
Résolvons $\displaystyle \frac{(x-1)(2x-3)}{x+1}>0$.

\textbf{C.E. :} $x+1\neq 0$, donc $x\neq -1$.

\textbf{Valeurs critiques :} $x=1$ (annule $x-1$), $x=\tfrac32$ (annule $2x-3$), $x=-1$ (interdite).
Ordonnées : $-1<1<\tfrac32$.

\medskip
\begin{center}
\renewcommand{\arraystretch}{1.45}
\setlength{\arraycolsep}{5pt}
\begin{tabular}{|>{\columncolor{primary!8}}l|c c c c c c c|}
\hline
\rowcolor{primary!15}
\textbf{Valeur de }$x$ & $-\infty$ & & $-1$ & & $1$ & $\tfrac32$ & $+\infty$ \\
\hline
Signe de $x-1$        & & $-$ & $|$ & $-$ & $0\ \ +$ & $+$ & \\
Signe de $2x-3$       & & $-$ & $|$ & $-$ & $-$ & $0\ \ +$ & \\
Signe de $x+1$        & & $-$ & $0$ & $+$ & $+$ & $+$ & \\
\hline
\rowcolor{excol!10}
Signe de $\dfrac{(x-1)(2x-3)}{x+1}$ & & $-$ & $\nexists$ & $+$ & $0\ \ -$ & $0\ \ +$ & \\
\hline
\end{tabular}
\end{center}

\medskip
\textbf{Lecture :} le quotient est strictement positif sur $]-1,1[$ et sur $]\tfrac32,+\infty[$. D'où
\[ \boxed{\,S=\,]-1,1[\ \cup\ \Bigl]\tfrac32,+\infty\Bigr[.\,} \]
\textbf{Attention au signe.} Si l'on multipliait les deux membres par un nombre \emph{négatif}, le sens
de l'inégalité changerait ; le tableau de signes évite ce piège en gardant $0$ à droite.
\end{exemple}

% =====================================================================
\section{Résolution des systèmes d'équations}
% =====================================================================
Un \textbf{système} est un ensemble d'équations qui doivent être satisfaites \emph{simultanément}.
Résoudre un système, c'est trouver les valeurs des inconnues qui conviennent à \emph{toutes} les
équations à la fois. Géométriquement, chaque équation linéaire représente une droite (à 2 inconnues)
ou un plan (à 3 inconnues) ; la solution est leur \emph{intersection}.

\subsection{Systèmes à deux inconnues}

Considérons le système
\[ \begin{cases} x+2y=5 & (\mathrm{I}) \\ 3x+y=20 & (\mathrm{II}) \end{cases} \]
Chaque équation est celle d'une droite ; la solution $(x,y)$ est le \textbf{point d'intersection} des
deux droites.

\begin{center}
\begin{tikzpicture}
  \begin{axis}[width=10.5cm,height=6cm,axis lines=middle,
      xmin=-1,xmax=9,ymin=-3,ymax=7,
      xlabel={$x$},ylabel={$y$},xtick={2,4,6,8},ytick={-2,2,4,6},grid=both,
      grid style={black!10},
      every axis x label/.style={at={(ticklabel cs:1)},anchor=west},
      every axis y label/.style={at={(ticklabel cs:1)},anchor=south},
      legend pos=south west,legend cell align=left,legend style={font=\scriptsize,draw=black!30}]
    \addplot[curveA,line width=1.2pt,domain=-1:9]{(5-x)/2};   % I : y=(5-x)/2
    \addplot[curveB,line width=1.2pt,domain=-1:9]{20-3*x};    % II : y=20-3x
    \addlegendentry{$(\mathrm{I})\ x+2y=5$}
    \addlegendentry{$(\mathrm{II})\ 3x+y=20$}
    \fill[black] (axis cs:7,-1) circle (2.2pt);
    \node[anchor=south west,font=\small] at (axis cs:7,-1){$(7,-1)$};
  \end{axis}
\end{tikzpicture}
\end{center}

\subsubsection{Méthode A — par substitution}

\begin{commentfaire}[substitution]
\begin{enumerate}[leftmargin=2em,topsep=2pt,itemsep=2pt]
  \item \textbf{Isoler} une inconnue dans l'une des équations.
  \item \textbf{Remplacer} cette inconnue par son expression dans l'\emph{autre} équation.
  \item \textbf{Résoudre} l'équation à une seule inconnue obtenue.
  \item \textbf{Réinjecter} la valeur trouvée pour calculer la seconde inconnue.
\end{enumerate}
\end{commentfaire}

\begin{exemple}[substitution, pas à pas]
\begin{enumerate}[leftmargin=2em,topsep=2pt,itemsep=2pt]
  \item On isole $x$ dans (I) : $\;x=5-2y$.
  \item On remplace dans (II) : $\;3(5-2y)+y=20$.
  \item On résout : $\;15-6y+y=20 \Rightarrow -5y=5 \Rightarrow y=-1$.
  \item On réinjecte : $\;x=5-2\times(-1)=7$.
\end{enumerate}
\[ S=\{(7\,;-1)\}. \]
\end{exemple}

\subsubsection{Méthode B — par combinaison linéaire}

\begin{commentfaire}[combinaison (addition)]
\begin{enumerate}[leftmargin=2em,topsep=2pt,itemsep=2pt]
  \item \textbf{Multiplier} chaque équation par un nombre bien choisi pour rendre \emph{opposés} les coefficients d'une inconnue.
  \item \textbf{Additionner} les deux équations : l'inconnue choisie disparaît.
  \item \textbf{Résoudre} l'équation à une inconnue restante.
  \item \textbf{Réinjecter} dans l'une des équations de départ.
\end{enumerate}
\end{commentfaire}

\begin{exemple}[combinaison, pas à pas]
On veut éliminer $x$. On multiplie (I) par $3$ et (II) par $-1$ pour obtenir des coefficients de $x$ opposés :
\[ \begin{cases} (x+2y)\times 3=5\times 3 \\ (3x+y)\times(-1)=20\times(-1)\end{cases}
   \;\Longleftrightarrow\;
   \begin{cases} +3x+6y=15 \\ -3x-y=-20\end{cases} \]
On additionne (I)$+$(II) : $\;6y-y=15+(-20)$, soit $5y=-5$, donc $y=-1$.
On réinjecte dans (I) : $\;x+2\times(-1)=5$, soit $x=7$.
\[ S=\{(7\,;-1)\}. \]
On retrouve bien la même solution : c'est le point d'intersection des deux droites.
\end{exemple}

\subsection{Systèmes à trois inconnues : le pivot de Gauss}

Pour trois inconnues, la méthode systématique est le \textbf{pivot de Gauss} : on élimine les inconnues
une à une pour rendre le système « triangulaire », puis on remonte.

\begin{definition}[méthode du pivot de Gauss]
On utilise une équation comme \textbf{pivot} pour éliminer une inconnue dans les autres équations, par
des opérations qui ne changent pas l'ensemble solution : remplacer une ligne $L_i$ par $L_i-k\,L_j$
(combinaison de lignes). On obtient un système \emph{échelonné} (triangulaire) que l'on résout de bas
en haut.
\end{definition}

\begin{commentfaire}[pivot de Gauss, à 3 inconnues]
\begin{enumerate}[leftmargin=2em,topsep=2pt,itemsep=2pt]
  \item \textbf{Pivot 1 :} garder $L_1$ ; s'en servir pour éliminer $x$ dans $L_2$ et $L_3$.
  \item \textbf{Pivot 2 :} garder $L_2$ (nouvelle) ; s'en servir pour éliminer une seconde inconnue dans $L_3$.
  \item \textbf{Triangle :} $L_3$ ne contient plus qu'une inconnue $\Rightarrow$ on la calcule directement.
  \item \textbf{Remontée :} réinjecter dans $L_2$, puis dans $L_1$, pour obtenir les autres inconnues.
\end{enumerate}
\end{commentfaire}

\begin{exemple}[résolution complète d'un système $3\times3$]
Résolvons
\[ \begin{cases} x+2y+z=4 & L_1 \\ 2x+y-z=5 & L_2 \\ 3x+y+2z=3 & L_3 \end{cases} \]

\textbf{Étape 1 — éliminer $x$.} $L_1$ est le pivot. On remplace $L_2\leftarrow L_2-2L_1$ et $L_3\leftarrow L_3-3L_1$ :
\[ \begin{cases} x+2y+z=4 & L_1 \\ -3y-3z=-3 & L_2 \\ -5y-z=-9 & L_3 \end{cases} \]

\textbf{Étape 2 — éliminer une inconnue de $L_3$.} On combine pour faire disparaître $z$ :
$L_3\leftarrow L_2-3L_3$ donne
\[ -3y-3z-3(-5y-z)=-3-3(-9)\;\Longrightarrow\;12y=24, \]
d'où le système triangulaire
\[ \begin{cases} x+2y+z=4 & L_1 \\ -3y-3z=-3 & L_2 \\ 12y=24 & L_3 \end{cases} \]

\textbf{Étape 3 — remontée.} De $L_3$ : $\;12y=24\Rightarrow y=2$. On reporte dans $L_2$ :
$-3(2)-3z=-3\Rightarrow -6-3z=-3\Rightarrow z=-1$. Enfin dans $L_1$ :
$x+2(2)+(-1)=4\Rightarrow x+3=4\Rightarrow x=1$.
\[ \boxed{\,S=\{(1\,;2\,;-1)\}.\,} \]
\end{exemple}

\subsection{Comment faire : mettre un problème concret en système}

Beaucoup d'énoncés « en français » se traduisent directement en systèmes. La compétence-clé est la
\emph{traduction} mot à mot.

\begin{commentfaire}[traduire un énoncé en équations]
\begin{enumerate}[leftmargin=2em,topsep=2pt,itemsep=2pt]
  \item \textbf{Nommer} les inconnues (par ex. $x,y,z$ = les prix recherchés).
  \item \textbf{Traduire} chaque phrase en une équation (« le triple de », « augmenté de », « la somme »\dots).
  \item \textbf{Nettoyer} : chasser les fractions (multiplier chaque ligne), ranger les termes.
  \item \textbf{Résoudre} par la méthode adaptée (Gauss pour 3 inconnues) et \textbf{interpréter} le résultat.
\end{enumerate}
\end{commentfaire}

\begin{exemple}[un problème de prix à trois inconnues]
\emph{Trois articles ont des prix $x$, $y$, $z$ (en \euro). On sait que :}
\begin{itemize}[leftmargin=1.6em,topsep=2pt]
  \item \emph{le prix du premier plus le triple du second et du troisième vaut \num{54}~\euro} : $\;x+3y+3z=54$ ;
  \item \emph{le double du premier plus le tiers du second vaut le troisième augmenté de \num{21}~\euro} : $\;2x+\tfrac{y}{3}=z+21$ ;
  \item \emph{le troisième plus le tiers des deux autres vaut \num{16}~\euro} : $\;z+\tfrac{x}{3}+\tfrac{y}{3}=16$.
\end{itemize}

\textbf{Étape 1 — chasser les fractions} (multiplier $L_2$ et $L_3$ par $3$) et ranger :
\[ \begin{cases} x+3y+3z=54 & L_1 \\ 6x+y-3z=63 & L_2 \\ x+y+3z=48 & L_3 \end{cases} \]

\textbf{Étape 2 — pivot de Gauss.} $L_2\leftarrow L_2-6L_1$ et $L_3\leftarrow L_3-L_1$ :
\[ \begin{cases} x+3y+3z=54 \\ -17y-21z=-261 \\ -2y=-6 \end{cases} \]
La dernière ligne donne directement $y=3$.

\textbf{Étape 3 — remontée.} Dans $-17y-21z=-261$ : $-17\times3-21z=-261\Rightarrow -21z=-210\Rightarrow z=10$.
Puis $L_1$ : $x+3\times3+3\times10=54\Rightarrow x+39=54\Rightarrow x=15$.
\[ S=\{(15\,;3\,;10)\}. \]
\textbf{Interprétation :} les articles coûtent respectivement \num{15}~\euro, \num{3}~\euro{} et \num{10}~\euro.
\end{exemple}

% =====================================================================
\section{Synthèse : la boîte à outils de l'algèbre}
% =====================================================================
Tous les outils de cette fiche se répondent. Le schéma ci-dessous résume comment ils s'articulent
autour d'un seul objectif : \emph{transformer un problème compliqué en problèmes simples}.

\begin{center}
\begin{tikzpicture}[every node/.style={font=\small},node distance=0.8cm]
  \node[rectangle,rounded corners=5pt,fill=primary,text=white,font=\bfseries,
        minimum width=4.2cm,minimum height=1cm,align=center] (top)
        {BOÎTE À OUTILS\\ALGÉBRIQUE};
  \node[rectangle,rounded corners=4pt,draw=defcol,fill=defcol!10,text width=3.0cm,align=center,
        minimum height=1.2cm,below left=1.4cm and 3.4cm of top] (a) {\textbf{Puissances \& racines}\\ \scriptsize exposants, $a^{m+n}$, $\sqrt[n]{\cdot}$};
  \node[rectangle,rounded corners=4pt,draw=thmcol,fill=thmcol!12,text width=3.0cm,align=center,
        minimum height=1.2cm,below left=1.4cm and 0.0cm of top] (b) {\textbf{Logarithmes}\\ \scriptsize réciproque de l'exponentielle};
  \node[rectangle,rounded corners=4pt,draw=formulacol,fill=formulacol!12,text width=3.0cm,align=center,
        minimum height=1.2cm,below right=1.4cm and 0.0cm of top] (c) {\textbf{Factorisation}\\ \scriptsize mise en évidence, Horner};
  \node[rectangle,rounded corners=4pt,draw=excol,fill=excol!12,text width=3.0cm,align=center,
        minimum height=1.2cm,below right=1.4cm and 3.4cm of top] (d) {\textbf{Systèmes}\\ \scriptsize substitution, Gauss};
  \draw[->,>=Stealth,thick,primary] (top) -- (a);
  \draw[->,>=Stealth,thick,primary] (top) -- (b);
  \draw[->,>=Stealth,thick,primary] (top) -- (c);
  \draw[->,>=Stealth,thick,primary] (top) -- (d);
\end{tikzpicture}
\end{center}

\begin{tcolorbox}[boxbase,colback=primary!4,colframe=primary,
    title={\textsc{Les réflexes à garder en tête}}]
\begin{itemize}[leftmargin=1.6em,topsep=2pt,itemsep=2pt]
  \item \textbf{Domaine d'abord} : avant tout calcul, vérifier les conditions d'existence (pas de division par $0$, pas de log d'un négatif, pas de racine d'un négatif).
  \item \textbf{Même base} : pour résoudre $a^{u}=a^{v}$, ramener à une base commune, puis égaler les exposants.
  \item \textbf{Le log fait descendre l'exposant} : $\ln(a^{u})=u\ln a$ — l'arme contre les inconnues en exposant.
  \item \textbf{Factoriser pour résoudre} : un produit nul $\Rightarrow$ au moins un facteur nul.
  \item \textbf{Discriminant} $\Delta=B^{2}-4AC$ : son signe donne le nombre de racines sans tout calculer.
  \item \textbf{Tableau de signes} : la façon sûre de résoudre une inéquation produit/quotient sans erreur de signe.
  \item \textbf{Systèmes} : substitution ou combinaison à 2 inconnues, pivot de Gauss à 3 inconnues.
\end{itemize}
\end{tcolorbox}

\vfill
\begin{center}
\small\color{black!55}
\rule{0.5\textwidth}{0.4pt}\\[4pt]
Fin de la Fiche 1 — \emph{Rappel d'outils et concepts de base en algèbre}.\\
Cours rédigé d'après l'extrait du livre ; figures originales en TikZ.
\end{center}

\end{document}
